\documentclass[11pt]{article}

\usepackage[margin=1in]{geometry}
\usepackage{graphicx}
\usepackage{amsmath}
\usepackage{amssymb}
\usepackage{amsthm}
\usepackage{booktabs}
\usepackage{longtable}
\usepackage{array}
\usepackage{hyperref}
\usepackage{xcolor}
\usepackage{enumitem}
\usepackage{microtype}
\usepackage{seqsplit}
\usepackage{float}
\usepackage[T1]{fontenc}
\usepackage[utf8]{inputenc}

\hypersetup{
  colorlinks=true,
  linkcolor=blue,
  citecolor=blue,
  urlcolor=blue
}

\theoremstyle{definition}
\newtheorem{definition}{Definition}[section]
\newtheorem{obstruction}[definition]{Obstruction}
\theoremstyle{plain}
\newtheorem{theorem}[definition]{Theorem}
\newtheorem{lemma}[definition]{Lemma}
\newtheorem{corollary}[definition]{Corollary}
\theoremstyle{remark}
\newtheorem{remark}[definition]{Remark}

\title{DASHI: A Constructive Carrier Derivation Program for Physics-Unification Structures\\
\large Ultrametric Carrier Geometry, Projection-Defect Semantics, and Mechanized Closure Frontiers}
\author{DASHI project}
\date{Draft of 2026-05-17}

\begin{document}
\sloppy
\maketitle

\begin{abstract}
Background: physics-unification programs often state target equations before
making explicit which carrier, quotient, source, and calibration obligations
are actually inhabited. Objective: this paper constructs a typed ultrametric
carrier geometry for staging those obligations without collapsing theorem,
target, obstruction, and empirical receipt into one status. Methods: DASHI
combines a ternary/refinement carrier, \texttt{FactorVec} valuation
coordinates, projection-defect decomposition, filtered operator surfaces, local
transport, root-shell boundaries, UFT/motif compression, and mechanized
admissibility receipts. Results: the manuscript derives a carrier spine from
unresolvedness through recursive refinement, ultrametric distance, tracked lane
commutation, projection residuals, and quotient-sensitive filtration; then
applies that spine to gauge curvature, Yang-Mills ordering defects, QFT
commutator targets, GR curvature/Bianchi targets, measurement residuals,
sector splitting, compression admissibility, and a P5' anomaly-candidate
empirical contact surface: \texttt{empiricalContactReached = true} for
\texttt{p5PrimeBorderlineAnomalyCandidate}, with residuals \(-2.8\sigma\)
and \(-3.0\sigma\) in the \(q^2\) bins \([4,6]\) and \([6,8]\)
\(\mathrm{GeV}^2/c^4\), under digest-bound SHA256 artifact discipline using
\texttt{flavio 2.7.0} and the BSZ baseline. The same typed discipline blocks
stronger readings: no new physics is claimed, no
\texttt{acceptedResidualCandidate} is inhabited, no full covariance matrix has
been found, upstream E8 completion is not promoted, and accepted empirical
authority, sourced non-flat GR, continuum recovery, GRQFT closure, and
completed unification remain open. Conclusion: DASHI proposes a constructive
ultrametric carrier geometry in which admissible refinement, compression,
transport, and quotient structure generate staged physics-facing derivation
surfaces, while typed closure semantics track the residual obligations that
block stronger promotion.
\end{abstract}

\noindent\textbf{Keywords:} ultrametric geometry; physics unification; gauge
theory; filtered operator algebras; projection-defect decomposition; dependent
type theory; constructive mathematics; formal methods.

\section{Introduction}

DASHI separates proved structure, target equations, empirical fits, and
unresolved obstructions inside one constructive carrier geometry. A state
carries trit-valued local data, tracked prime-lane valuations, ultrametric refinement
depth, lane-local actions, and projections into observable surfaces. The
residual left by a projection remains part of the formal state rather than
being discarded or promoted by convention.

Observable structure is modeled as a projection from a richer carrier space.
For a surface \(s\), a projection \(P_s:X\to Y_s\) exposes the admissible
component of \(x\in X\). The defect \(D_s(x)\) records the residual structure:
unresolved branches, incompatible quotient data, missing adapters, or absent
empirical authority. Filtrations organize limiting behavior and associated
graded constructions. Ultrametric refinement organizes address depth and
precision. Independent lane actions give componentwise laws for selected
coordinate sectors.

The same carrier supports local transport, filtered operators, ultrametric
refinement, finite root-shell boundaries, compression structures, and bounded
empirical receipts. Direction-indexed discrete sites define local transport
schemas. Finite-support operators define filtered contraction targets. The
tracked G6 commuting theorem gives a coordinate-independence law for selected
prime-lane actions. Local E8/LILA surfaces record bounded root-geometry and
upstream-promotion boundaries. UFT-style compression records repeated
refinement motifs and recoverable residual distinctions. The primary empirical
contact in this draft is the P5' anomaly candidate: the receipt records
\texttt{empiricalContactReached = true} for
\texttt{p5PrimeBorderlineAnomalyCandidate}, with \(-2.8\sigma\) and
\(-3.0\sigma\) residuals in the \(q^2\) bins \([4,6]\) and \([6,8]\)
\(\mathrm{GeV}^2/c^4\), bound to SHA256-digested artifacts and a frozen
\texttt{flavio 2.7.0} + BSZ baseline. Drell-Yan remains a methodology
demonstration for typed empirical receipts, not the lead empirical claim.

\section*{Target Obligation Surfaces and DASHI Derivation Roles}

The governing order for every displayed equation is
\[
\text{carrier schema} < \text{interpretation map} <
\text{field-equation law} < \text{calibrated physics claim}.
\]
Until the interpretation map, quotient laws, physical units, source interface,
and empirical/authority receipts are inhabited, the equations in this opening
section are target obligation surfaces. They are placed here first because they
are the reason the derivation spine matters, not because they are already
closed.
DASHI's current claim is that it gives a typed derivation route in which each
term has an explicit carrier-side obligation:
\[
\Omega_p=H_p-I,
\qquad
F_A=dA+A\wedge A,
\qquad
[\nabla_\mu,\nabla_\nu]V^\rho
=R^\rho{}_{\sigma\mu\nu}V^\sigma,
\]
\[
R_{\mu\nu}=R^\rho{}_{\mu\rho\nu},
\qquad
G_{\mu\nu}=R_{\mu\nu}-\frac12Rg_{\mu\nu},
\]
\[
G_{\mu\nu}+\Lambda g_{\mu\nu}
=
\frac{8\pi G}{c^4}T_{\mu\nu}.
\]
The gauge and operator targets are similarly explicit:
\[
d_AF_A=0,
\qquad
d_A{*F_A}=J,
\qquad
dF=0,\quad d{*F}=J,
\]
\[
U(t)=e^{-itH},
\qquad
i\partial_t\psi=H\psi,
\qquad
[\phi(t,x),\pi(t,y)]=i\delta(x-y).
\]

\begin{center}
\small
\begin{tabular}{p{0.22\linewidth}p{0.37\linewidth}p{0.32\linewidth}}
\toprule
Physics term & DASHI structural reading & First residual obligation \\
\midrule
\(g_{\mu\nu}\) & projected metric over refinement carrier & metric emergence, inverse, signature, units \\
\(\nabla\) & admissible transport law & non-flat connection/transport adapter \\
\(R^\rho{}_{\sigma\mu\nu}\) & noncommuting refinement transport defect & curvature carrier, covariance, torsion convention \\
\(R_{\mu\nu},R\) & contracted transport defect & Ricci/scalar contraction and index laws \\
\(G_{\mu\nu}\) & stable curvature-obstruction tensor & contracted Bianchi/divergence-free law \\
\(T_{\mu\nu}\) & source/stress-energy target projection / matter residual carrier & stress-energy interface and authority \\
\(\Lambda g_{\mu\nu}\) & background refinement/vacuum term & cosmological convention and calibration \\
\(8\pi G/c^4\) & source-coupling normalization & weak-field/Newtonian limit and unit calibration \\
\bottomrule
\end{tabular}
\end{center}

The normalization \(\kappa=8\pi G/c^4\) is not arbitrary decoration. In a
completed route it must be fixed by geometric normalization, spherical-flux
structure, and weak-field consistency with the Newtonian limit. Paper 1 does
not prove that calibration. It makes the missing calibration a named
obligation rather than hiding it behind the formal carrier.

\section*{Core Carrier Grammar}

\noindent\fcolorbox{black!45}{black!3}{%
\begin{minipage}{\dimexpr\linewidth-2\fboxsep-2\fboxrule\relax}
\small
The constructions in this paper use a carrier of the form
\[
x \in \mathcal X \simeq (a,\nu,\tau,d,\ell,r),
\qquad
\tau \in \{-1,0,+1\}.
\]
Here \(a\) is a refinement address, \(\nu\) is
\texttt{FactorVec}-style prime-lane valuation data, \(\tau\) is the ternary
coordinate, \(d\) is refinement depth, \(\ell\) is the active lane, and \(r\)
records residual/projection status. A surface \(s\) exposes only its admitted
projection:
\[
x = P_s(x) + D_s(x), \qquad
\rho(x,y)=\alpha^{k(x,y)},\; 0<\alpha<1.
\]
\(P_s(x)\) is the promoted component; \(D_s(x)\) is the residual defect; and
\(k(x,y)\) is the first incompatible address, valuation, or prefix depth.
Filtered operator surfaces use \(F_0 \subseteq F_1 \subseteq \cdots\),
\(\operatorname{gr}_nF=F_n/F_{n-1}\), and descended brackets only where
quotient, equivalence, norm, and bracket laws are inhabited. Lane actions
commute only on the official tracked independent \texttt{GL.FactorVec} route.
UFT/motif compression records reusable projected structure while residual
defects remain visible.
\end{minipage}}

\vspace{0.8em}

\section*{Derivation Map at a Glance}

The closure frontier appears later. The early reader map is a derivation-role
map: each target equation is read as a structural role to be produced by the
carrier, followed by the adapter still needed for a calibrated physics claim.
\begin{center}
\small
\begin{tabular}{p{0.20\linewidth}p{0.38\linewidth}p{0.32\linewidth}}
\toprule
Target role & DASHI derivation engine & Adapter still needed \\
\midrule
Maxwell / abelian gauge & commuting local transports give additive plaquette curvature \(F=dA\) with \(dF=0\) from \(d_1d_0=0\) & bind adopted G2 schema into Maxwell scope; exterior derivative/nilpotency, Hodge star, source current, U(1) sector/action, covariance, units/convention \\
Yang-Mills / nonabelian gauge & noncommuting local transports force the commutator term \(A\wedge A\) in \(F_A=dA+A\wedge A\) & representation, sector restriction, action/source law, calibrated field equation and coupling convention \\
QFT commutator & filtered finite-support operators descend only through quotient-correct brackets & equivalence modulo denominator filtration, accepted quotient carrier, descended norm/bracket, Hilbert/amplitude/CCR/Born laws \\
GR curvature & path-dependent carrier transport yields curvature as the minimal commutator obstruction & non-flat connection, Christoffel-from-metric law, Ricci/Bianchi reduction, stress-energy interface, units/calibration/source law \\
Measurement / uncertainty & projection \(\pi:X\to Q\) leaves residual fibres \(R_\pi(x)\) & amplitude/probability law plus empirical projection and observable map \\
Sector splitting & projection preserves a subgroup \(H\subseteq G\), leaving residual cosets \(G/H\) & concrete SM representation, subgroup projection law, coupling/mixing calibration \\
\bottomrule
\end{tabular}
\end{center}


\section{Background and Related Work}

DASHI draws on standard mathematical neighborhoods including ultrametric and
p-adic structure \cite{schikhof-1984,vladimirov-volovich-zelenov-1994},
lattice and gauge language
\cite{nakahara-2003,wilson-1974,kogut-1979}, contraction and associated
graded constructions \cite{weibel-1994}, root-system classification
\cite{humphreys-1972},
mechanized mathematics \cite{norell-2007,hott-book,plfa-agda}, and
compression-style trie geometry \cite{rissanen-1978,grunwald-2007}. The
claims made here are restricted to the repo-local receipt and obstruction
surfaces cited in the claim ledger.

Agda is used as the mechanized verification surface for selected propositions,
receipts, audits, and blockers. A proved construction is represented by an
inhabited artifact; a missing constructor is recorded as a manuscript-relevant
obligation. This places the paper near proof assistants, dependent type theory,
mechanized mathematics, and proof-carrying artifacts, while remaining narrower
than an end-to-end certified physics pipeline.

Ultrametric and p-adic language organizes refinement by depth: nearby states
share high-level valuation or prefix structure, while separation can occur at
the first incompatible digit or coordinate. Hensel-style lifting is used only
as an external analogy for compatible refinement across precision levels. In
Paper 1 this is expository support for carrier and compression geometry, not
an independent proof of convergence.

The labels G2, G3, and G6 are DASHI derivation-lane names. They should not be
read as claims of equivalence to standard mathematical objects with similar
names, such as the exceptional Lie group \(G_2\), unless an inhabited surface
states that equivalence. E8 is different: the paper explicitly uses standard
root-system language there, but only for the bounded local enumeration and
promotion-boundary surfaces stated below.

\subsection*{External Formal Bridges and Terminal Boundaries}

DASHI is not presented as an isolated replacement for existing formal
geometry or algebraic QFT. It is a carrier-and-receipt layer that can be
bridged to external formal mathematics when the corresponding adapters are
made explicit. The current bridge layer has six named obligation slots,
B0--B5, each recorded as a typed obligation rather than as a prose promise.

\textbf{B0 (geometric emergence).} Wave-coherent and refinement-stable DASHI
transport should determine a DCHoTT \texttt{G-structures} socket. This is the
central open bridge theorem and the target of the next paper. The current
receipt is \texttt{LeviCivitaBridge.agda}: it typechecks against the cloned
DCHoTT-Agda flat modules and records B0 as a named postulated obligation, not
an imported theorem \cite{wellen-2018,dchott-agda}.
\texttt{DCHoTTBridgeObligationIndex.agda} further names the B0.1 pro-object
compatible-family/formal-disk scaffold, the B0.2 flat-formal-disk target, and
the B0.3 refinement-stable frame/metric/G-structure reduction target.
\texttt{ProObjectSemantics.agda} constructs the DASHI-side B0.1
compatible-family limit surface: pro-object points, depth projections,
refinement coherence, depthwise agreement balls, and the depth-zero
formal-thickening relation. The analytic and external-geometric steps remain
open: real-valued ultrametric completion, transport-smooth sheaves,
formal-disk-to-DCHoTT equivalence, DCHoTT manifold/G-structure promotion, and
Levi-Civita specialization are not proved here.

\textbf{B1 (finite curvature and continuum boundary).} The curvature surface
is no longer only a nilpotent-transport target. The repository has a finite
full depth-9 connection-and-curvature receipt and a finite Yang--Mills field
equation surface for the source-free equation \(D\star F=0\). These are
finite carrier-level receipts for curvature, Bianchi/Yang--Mills algebra, and
field-equation shape; they do not import continuum physics by themselves. The
non-flat metric/Levi-Civita, stress-energy, Wald/continuum GR authority, and
continuum promotion steps remain external bridge boundaries downstream of B0.

\textbf{B2 (local quantum field).} The cited Haag-Kastler stack route is
recorded as identifying a locally covariant AQFT as a point of the relevant HK
2-functor. The finite prime-lane DHR surface has also advanced: local
endomorphism/category receipts, braiding naturality, left and right hexagon
laws, statistics-as-braiding, and conjugate zigzags are finite-level inhabited
receipts. The Klein-Gordon free-field witness is recorded in
\texttt{KleinGordonAQFTReceipt.agda}, and
\texttt{AQFTCarrierAlgebraQuotientSurface.agda} stages \(A(O)\), \(A_d(O)\),
the filtered colimit \(\mathrm{colim}_d A_d(O)\), and Cauchy time-slice
evolution as typed targets only. The arbitrary DHR symmetric monoidal
structure, Doplicher--Roberts theorem application, compact group
reconstruction, and exact \(G_{\mathrm{DHR}}\cong G_{\mathrm{SM}}\) comparison
remain blocked or external \cite{haag-kastler-stacks-2024}.

\textbf{B3 (interacting boundary).} Interacting QFT, renormalisation, and
coupling calibration are typed adapter boundaries. The record
\texttt{InteractingQFTBoundaryReceipt.agda} treats constructive interacting
local algebra nets as the first missing primitive and explicitly blocks
Standard Model or interacting-QFT promotion.

\textbf{B4 (compression admissibility).} Compression is governed by an MDL
Lyapunov condition: admissible transport must not increase the required
description length at the active projection grade. This is recorded by the
compression admissibility receipt and used in the compression section below.

\textbf{B5 (honest frontier).} The terminal status layer separates weak
modulo accounting from unqualified closure. The repository records a weak
terminal claim only modulo an explicit minimal package: B0.1 \(\Im\)
reflection, B0.3 \texttt{WeakBG}, Bisognano--Wichmann/time-slice authority,
Doplicher--Roberts reconstruction, \(C^*\)-completion, and external/Clay
mass-gap authority. DASHI-internal ST3/Banach mass-gap closure is recorded
separately and does not promote Clay authority. This is not a terminal theorem.
The unqualified terminal claim remains blocked; Clay/external mass-gap
authority, DHR-to-Standard-Model matching, final adapter acceptance,
stress-energy/Wald authority, and GRQFT composition remain open gates.
\texttt{laneDimension} has only a conditional internal Moonshine bridge route:
Ogg genus-zero selects the 15 supersingular primes, no-ghost supplies orbifold
compatibility, Dong-Li-Mason supplies the weight-1 Lie-algebra table authority,
and \texttt{DASHIPrimeLaneEquiv} must bridge that table to DASHI lane
dimensions. The repository files
\texttt{ExternalFormalImportRoadmapReceipt.agda},
\texttt{BalabanRGMassGapReceiptSurface.agda},
\texttt{TerminalOpenProblemStatusSurface.agda}, and
\texttt{GRQFTTerminalCompositionBoundary.agda} record these terminal boundaries
without promoting completed unification.

The downstream Papers 2--7 map is a named obligation ledger. Its current core
inventory is organized into eight gates: carrier-to-physics functor laws, S8
spectral-gap witnesses, nonabelian Yang--Mills structure, non-flat FactorVec
GR, Hilbert/Stone depth limits, DHR/AQFT reconstruction, matter/Higgs/Yukawa
and CKM arithmetic, and final Paper7Receipt composition. Recent repository
work has typed target surfaces for these gates, including a Nat S8 witness
slice, finite \texttt{SU(2)} Jacobi support, a p2-time valuation metric
candidate, finite-support depth-colimit targets, and supplied-sector DHR
tensor/fusion sockets. These are local proof surfaces, not completed gate
receipts.

The active core theorem roadmap is deliberately smaller than the application
surface inventory. Chemistry, biology, cognition, rendering, culture, and other
application lanes remain downstream/frozen unless they directly discharge one
of the eight core gates. The map does not alter
\texttt{terminalClaimPromoted = false},
\texttt{continuumClayMassGapPromoted = false}, or any external-authority
boundary.

The route manifest separates similarly named blockers. Route A
\texttt{d\^{}2=0} is an abelian exterior-derivative nilpotency route from
prime-coordinate commutativity / boundary-of-boundary cancellation; the current
mechanized surface also tracks coefficient laws and SFGC action-bridge
requests. Route B nonabelian Yang-Mills keeps the Killing-form Ad-invariance,
Hodge/covariant derivative, variation/IBP, and current/source residuals
visible. The GR boundary is separate again: finite Bianchi, CRT, or sidecar
surfaces do not supply non-flat Levi-Civita, Ricci/Einstein, stress-energy, or
continuum GR promotion. Gate 1 now has a graph-level carrier-to-physics
interpretation map for p2/p3/p5/p7, but physical covariance and
filtered-colimit preservation remain open.

The honest frontier claim is therefore deliberately weaker than a completed
unification theorem: DASHI localizes which pieces are derived from carrier
transport and which pieces require external geometric, algebraic, empirical,
or calibration input. The absence of an inhabitant is manuscript-relevant
state, not an editorial disclaimer.


\section{Derivation Spine}

The carrier geometry is not introduced as an inventory. It is derived from a
small sequence of representational obligations: unresolved local state,
coordinate traversal, recursive refinement, prime-lane valuation, surface
projection, residual defect, and filtered operator descent.

\begin{center}
\begin{tabular}{p{0.16\linewidth}p{0.28\linewidth}p{0.46\linewidth}}
\toprule
Stage & Object & Spine statement \\
\midrule
1 & Primitive ternary state & Signed evidence with unresolved residue requires at least three local statuses. \\
2 & Traversal & Coordinate-local updates turn primitive states into paths in a finite state graph. \\
3 & Voxel / hypercube cell & A finite block of ternary coordinates is a \(3^k\) cell with coordinate-wall adjacency. \\
4 & Trie and ultrametric & Recursive refinement is represented by prefix addresses whose natural distance is ultrametric. \\
5 & \texttt{FactorVec} & Prime-lane valuation counts refine addresses into the tracked \texttt{Vec15 Nat} carrier. \\
6 & Projection-defect & Every visible reading is a projection together with a residual defect. \\
7 & Filtration & Depth, support, valuation, and admissibility organize surfaces into filtered theorem targets. \\
\bottomrule
\end{tabular}
\end{center}

\begin{definition}[Primitive local state]
Let \(T=\{-1,0,+1\}\). A primitive local state is an element \(t\in T\).
The values \(+1\) and \(-1\) record oriented alternatives. The value \(0\)
records unresolved residue: structure that has not been discharged by the
current surface.
\end{definition}

\begin{lemma}[Ternary minimality]
Any local carrier that distinguishes a positive orientation, a negative
orientation, and unresolved residue needs at least three local values.
\end{lemma}
\begin{proof}
A two-valued carrier distinguishes at most two equivalence classes. If
unresolved residue is identified with \(+1\), unresolved structure is promoted
as positive evidence. If it is identified with \(-1\), unresolved structure is
promoted as negative evidence. If \(+1\) and \(-1\) are identified, orientation
is lost. The three roles therefore require three values. This is a bookkeeping
minimality result, not a physical law.
\end{proof}

\begin{definition}[Traversal]
For a finite coordinate set \(I\), a primitive state is a function \(x:I\to T\).
A coordinate update is a map \(u_{i,\phi}:T^I\to T^I\) that changes only
coordinate \(i\), using a local rule \(\phi:T\to T\), and leaves every
\(j\neq i\) fixed. A traversal is a finite sequence
\[
x_0 \xrightarrow{u_1} x_1 \xrightarrow{u_2}\cdots
\xrightarrow{u_n} x_n .
\]
\end{definition}

\begin{lemma}[Traversal as address word]
Every traversal determines a word in the update alphabet
\(\Sigma=\{(i,\phi)\mid i\in I,\phi:T\to T\}\).
\end{lemma}
\begin{proof}
Record the label \((i_m,\phi_m)\) of each update \(u_m\). The traversal is
represented by the word
\[
(i_1,\phi_1)(i_2,\phi_2)\cdots(i_n,\phi_n)\in\Sigma^\ast.
\]
The word is a formal update trace relative to the chosen carrier; it does not
by itself claim a unique physical history.
\end{proof}

\begin{definition}[Ternary cell and voxel refinement]
For a finite block \(B\subseteq I\) with \(|B|=k\), the local cell is
\(C_B=T^B\). It has \(3^k\) vertices. Two vertices are coordinate-adjacent
when they differ in exactly one coordinate. A voxel refinement replaces one
coordinate \(i\) by a finite child block \(B_i\), replacing the factor \(T\) at
\(i\) by \(T^{B_i}\).
\end{definition}

\begin{lemma}[Hypercube language is finite-coordinate language]
The cube, voxel, and hypercube terms used in the origins story can be read as
finite coordinate cells and their refinements.
\end{lemma}
\begin{proof}
A block \(T^B\) is a finite product of local coordinate carriers. Its
coordinate-adjacency graph is the ternary analogue of a cubical graph: faces
are obtained by fixing coordinates, and refinements replace a coordinate by a
child product. Thus the geometric language refers to finite-coordinate
movement and refinement, not to an assertion that physical spacetime has been
derived.
\end{proof}

\begin{definition}[Refinement trie and ultrametric]
Let \(\Sigma\) be a finite refinement alphabet. A trie address is a finite or
infinite word in \(\Sigma\). For infinite addresses \(a,b\), let
\(\lambda(a,b)\) be the length of their longest common prefix, with
\(\lambda(a,a)=\infty\). Fix \(0<\alpha<1\), and define
\[
\rho(a,b)=
\begin{cases}
0, & a=b,\\
\alpha^{\lambda(a,b)}, & a\neq b.
\end{cases}
\]
\end{definition}

\begin{lemma}[Trie distance is ultrametric]
The distance \(\rho\) satisfies
\[
\rho(a,c)\leq \max(\rho(a,b),\rho(b,c)).
\]
\end{lemma}
\begin{proof}
If \(a\) and \(b\) agree for \(m\) symbols, and \(b\) and \(c\) agree for
\(m\) symbols, then \(a\) and \(c\) also agree for \(m\) symbols. Hence
\(\lambda(a,c)\geq \min(\lambda(a,b),\lambda(b,c))\). Since
\(0<\alpha<1\), increasing prefix depth decreases distance, so the displayed
ultrametric inequality follows.
\end{proof}

\begin{definition}[\texttt{FactorVec}]
Let
\[
P_{\mathrm{SS}}=\{2,3,5,7,11,13,17,19,23,29,31,41,47,59,71\}
\]
be the tracked supersingular-prime lane set used by the repo-native carrier.
The set \(P_{\mathrm{SS}}\) is characterized as the supersingular prime set in
Ogg's 1975 genus-zero observation. The repo receipt
\texttt{DASHI/Physics/Moonshine/SupersingularPrimeLaneBridge.agda} records this
bridge and explicitly keeps
\texttt{primeSetForcedFromFirstPrinciples = false}: the choice is a motivated
DASHI design decision, not a first-principles derivation.
A \texttt{FactorVec} is a function \(\nu:P_{\mathrm{SS}}\to\mathbb{N}\), equivalently a
\texttt{Vec15 Nat} exponent vector with one coordinate for each tracked prime
lane. For \(p\in P_{\mathrm{SS}}\), define the prime-lane bump \(B_p\) by
\[
B_p(\nu)(q)=
\begin{cases}
\nu(q)+1, & q=p,\\
\nu(q), & q\neq p.
\end{cases}
\]
\end{definition}

\begin{definition}[Supersingular-prime lane bridge]
The tracked lane set \(P_{\mathrm{SS}}\) is not an arbitrary list. It is the 15
supersingular-prime set
\[
\mathcal{P}_{SS}
=\{2,3,5,7,11,13,17,19,23,29,31,41,47,59,71\}.
\]
The receipt
\texttt{DASHI/Physics/Moonshine/SupersingularPrimeLaneBridge.agda} records the
external Ogg authority that the following three characterizations coincide:
\(X_0^+(p)=X_0(p)/w_p\) has genus zero; the supersingular polynomial in
characteristic \(p\) splits over \(\mathbb{F}_p\); and \(p\) divides the order
of the Monster group \cite{ogg-1975,borcherds-1992}. DASHI uses the second
condition as the carrier-facing entry point: a tracked lane is field-complete
at local depth \(1\) when all supersingular \(j\)-invariants are already
defined over \(\mathbb{F}_p\).
\end{definition}

\begin{remark}[Governance status of the 15SSP bridge]
The bridge is deliberately weaker than a first-principles derivation. The
receipt inhabits \texttt{DASHIPrimeSetIsP\_SS = true}, while keeping
\texttt{primeSetForcedFromFirstPrinciples = false},
\texttt{oggOriginalQuestionResolved = false}, and
\texttt{standardModelGaugeGroupDerived = false}. Thus the paper may claim that
the prime set has a precise Ogg-supersingular characterization, but not that
DASHI derives that set uniquely from the empty carrier or derives the Standard
Model gauge group from Moonshine.
\end{remark}

\begin{lemma}[Coordinate bumps commute]
For all \(p,q\in P_{\mathrm{SS}}\), \(B_p(B_q(\nu))=B_q(B_p(\nu))\).
\end{lemma}
\begin{proof}
Check each coordinate \(r\in P_{\mathrm{SS}}\). If \(r\) is neither \(p\) nor \(q\), both
sides equal \(\nu(r)\). If \(r=p\neq q\), both sides equal \(\nu(p)+1\). If
\(r=q\neq p\), both sides equal \(\nu(q)+1\). If \(p=q=r\), both sides equal
\(\nu(p)+2\). The vectors are equal coordinatewise.
\end{proof}

\begin{definition}[Surface projection and defect]
A surface \(s\) consists of a visible carrier \(Y_s\), an admissibility gate
\(G_s\), and a projection \(P_s:X\to Y_s\). The defect \(D_s(x)\) is the
residual information not promoted by \(P_s\). When the carrier has an additive
split, this is written \(x=P_s(x)+D_s(x)\).
\end{definition}

\begin{lemma}[Defect makes non-closure visible]
If two states \(x,x'\in X\) satisfy \(P_s(x)=P_s(x')\) but differ in gate
status, residual content, or unmet obligations, then the projection alone
cannot distinguish them. Recording \((P_s(x),D_s(x))\) keeps the promoted
component and the unpromoted residual in the same formal account.
\end{lemma}

\begin{definition}[Filtration and associated graded target]
A filtration on \(X\) is an increasing sequence
\[
F_0\subseteq F_1\subseteq F_2\subseteq\cdots .
\]
When quotient structure is available, the \(n\)-th graded piece is
\(\operatorname{gr}_nF=F_n/F_{n-1}\). This notation is a target unless the
equivalence relation, quotient carrier, projection-respects-equivalence law,
and descended operations are supplied.
\end{definition}

\begin{lemma}[Descent requires equivalence preservation]
Suppose \(A:X\to X\) is filtration-preserving and \(x\sim_n y\) means
\(x-y\in F_{n-1}\). If \(x\sim_n y \Rightarrow A(x)\sim_n A(y)\), then \(A\)
induces an operator on \(F_n/F_{n-1}\).
\end{lemma}
\begin{proof}
Define the induced map by \([x]\mapsto[A(x)]\). If \([x]=[y]\), then
\(x\sim_n y\). By equivalence preservation, \(A(x)\sim_n A(y)\), so
\([A(x)]=[A(y)]\). Thus the definition is independent of representative.
\end{proof}

\begin{theorem}[Early DASHI carrier spine]
Starting from primitive signed state with unresolved residue and the chosen
supersingular-prime lane set, the carrier grammar follows the staged design
spine recorded here. The route is constrained by the preceding obligations, but
it is not claimed as a first-principles forcing theorem:
\[
\begin{aligned}
\text{primitive state}
&\to \text{traversal}
\to \text{voxel / hypercube cell}\\
&\to \text{trie / ultrametric refinement}
\to \text{\texttt{FactorVec} valuation}\\
&\to \text{projection-defect}
\to \text{filtration}.
\end{aligned}
\]
\end{theorem}
\begin{proof}
The lemmas above derive the route. Ternary minimality gives primitive local
state. Traversal turns such states into coordinate-local update words. Ternary
cells give the bounded reading of voxel and hypercube traversal. Recursive
refinement gives prefix geometry and hence an ultrametric. The selected
prime-lane coordinate set is now governed by the 15SSP bridge: it is the
Ogg-supersingular set, with field-completeness over \(\mathbb{F}_p\) as the
DASHI-natural condition. Prime-lane refinement then gives the
\texttt{FactorVec} valuation route and basic coordinate commutation. Surface
observation motivates projection-defect accounting. Filtered
operator language then requires quotient descent before theorem-facing
dynamics can be promoted.
\end{proof}

The prime set itself is not derived from first principles in this theorem.
\texttt{DASHI/Physics/Moonshine/SupersingularPrimeLaneBridge.agda} records
\(\mathcal{P}_{SS}\) as the supersingular-prime set characterized by Ogg's
1975 classification, with Borcherds' Moonshine theorem as the external
Moonshine authority. The receipt keeps
\texttt{primeSetForcedFromFirstPrinciples = false}: the choice of
\(\mathcal{P}_{SS}\) is a motivated design decision, while the theorem above
states only what follows once that lane set is selected.

This theorem does not assert that the grammar completes physics or that
\(\mathcal{P}_{SS}\) is forced from first principles. It asserts that the
listed structures are not a catalog: they arise as successive obligations once
signed unresolved state, traversal, recursive refinement, the externally
characterized supersingular-prime lane set, surface projection, and filtered
operator descent are all kept visible.

\section{Core Constructive Architecture}

The carrier used below has six recurring components:
\[
  x \in X \simeq (a,\nu,\tau,d,\ell,r),
\]
where \(a\) is a refinement address, \(\nu\) is tracked prime-lane valuation
data, \(\tau\) is local trit state, \(d\) is depth or filtration grade, \(\ell\)
records the active lane/action context, and \(r\) records residual structure.
The later gauge, operator, root-shell, compression, and empirical sections use
different projections of this same carrier.

\subsection{Carrier Coordinates and Depth}

States are indexed by ternary local coordinates, refinement depth, and tracked
prime-lane valuation vectors. Prime and supersingular-prime coordinates provide
addressable lanes. Trit and Base369-style carriers distinguish positive,
negative, and unresolved positions. A \texttt{FactorVec} assigns one valuation
coordinate to each tracked prime lane, so lane updates can be checked
componentwise.

Trie addresses induce an ultrametric: states are close when they share a long
common refinement prefix, and they separate at the first incompatible depth.
The carrier therefore records depth and precision without treating every
refinement branch as an observed physical fact. Hensel-style lifting remains an
external analogy for compatible partial information across increasing
precision.

\subsection{Projection-Defect Decomposition}

The central local equation is the projection-defect split:
\[
  x = P(x) + D(x).
\]
Here \(P(x)\) denotes the promoted projection visible to a chosen surface, and
\(D(x)\) denotes the residual defect left by that projection. Projection exposes
the admissible component of a carrier state while retaining the unresolved
residual. Later promotion requires a typed bridge from \(D(x)\) into the
stronger surface.

\subsection{Filtration, Lane Action, and Compression}

Filtrations organize limiting behavior and associated graded constructions.
For an operator surface, the key question is whether restricted operators
descend to \(F_n/F_{n-1}\), whether the quotient is well defined, and whether
norms or brackets are preserved after projection. Lane actions are
coordinate-local transformations on tracked sectors; the tracked G6 commuting
theorem checks selected cross-lane actions on the current prime-valuation
route.

Repeated local substructures permit refinement-address compression without
collapsing admissible distinctions. UFT addresses record refinement structure,
motifs record reusable projected substructure at a chosen grade, and residuals
record distinctions that compression must preserve. A motif is not equality; it
is evidence of reusable projection structure subject to the compression
admissibility theorem below.

\subsection{Closure Discipline}

Closure in DASHI is a mathematical typing rule, not an editorial disclaimer. A
surface promotes only the projection whose evidence is inhabited at that
surface. The residual part remains visible as a defect: a missing quotient
law, adapter, authority token, convention receipt, calibration, or upstream
constructor.

This gives the manuscript one discipline for all lanes. A theorem surface may
promote a local algebraic statement; an empirical receipt may promote a
bounded comparison; a boundary record may promote only the fact that a
stronger bridge is still absent. Diagnostics, analogies, provenance, and
request surfaces can motivate later work, but they do not promote stronger
physical claims.

\begin{quote}
\texttt{surface S promotes projection P\_s(x);}\\
\texttt{residual D\_s(x) records the first missing bridge for stronger claim Q.}
\end{quote}

The intervening sections use this rule locally: each lane states its promoted
surface, its first residual obligation, and no stronger reading unless the
corresponding surface is inhabited.

\section{Mathematical Constructions as Applications}

The G2, G3, G6, and E8/LILA constructions instantiate the derivation spine on
specific local carriers: discrete plaquette transport, finite-support
operators, prime-lane commuting updates, and bounded root-shell geometry.

\begin{longtable}{p{0.12\linewidth}p{0.25\linewidth}p{0.28\linewidth}p{0.24\linewidth}}
\toprule
Label & Mathematical role & Current positive surface & Boundary \\
\midrule
G2 & Directed plaquette carrier schema & Direction-indexed SFGC schema & No Maxwell closure or field equation follows \\
G3 & Finite-support operator filtration & Selected finite-support subtraction support & Associated-graded quotient descent remains open \\
G6 & Prime-lane coordinate independence & Tracked cross-lane commuting theorem & Universal lane-operator promotion remains blocked \\
E8/LILA & Finite root-shell enumeration & Integer/two-sparse local completeness & Upstream E8 completion remains blocked \\
\bottomrule
\end{longtable}

\subsection{G2: Local Plaquettes from Coordinate Traversal}

The local carrier is modeled by directed plaquette neighborhoods. A site has
two visible movement directions; the minimal local cell records traversal in
one coordinate, traversal in the other, and comparison around the resulting
small square. This defines a checked interface for sites, directed edges, and
plaquette-shaped neighborhoods.

The positive surface is the G2 schema-extension receipt together with the G2
no-remaining-missing witness. Direction-indexed SFGC schema adoption is complete
for the current core surfaces. Dynamics remain open: curvature, field
equations, continuum recovery, and calibration are deferred to the closure
frontier. Exact Agda identifiers are listed in the reproducibility table.

\subsection{G3: Selected Finite-Support Operators and the Quotient Gap}

The G3 application starts where local traversal becomes operator calculus. The
current constructive result is certified subtraction support for a selected
constructive finite-support subtype, not arbitrary operator closure. Subtraction
closure is required before filtered operators descend to stable quotient
structures. The relevant inhabited surfaces are the selected finite-support
subtraction support receipt, the selected subtraction support witness, and the
selected subtraction certified-support record. A genuine associated graded layer still
requires quotient classes of the form \(F_n/F_{n-1}\) and the kernel-equivalence
law:
\[
  [x] = [0] \quad \Longleftrightarrow \quad
  x \textrm{ lies in the denominator filtration piece.}
\]
Until the accepted equivalence, quotient carrier,
projection-respects-equivalence, and descended norm/bracket/isomorphism laws
are inhabited, quotient notation is only a target shape. The current positive
surface is selected finite-support subtraction support. The open obstruction is
quotient-correct norm, product, bracket, and contraction structure.

\subsection{G6: Distinct-Coordinate Commutation in Tracked \texttt{FactorVec}}

The G6 lane records cross-lane independence inside the prime-coordinate
carrier. An update to the \texttt{p} lane and an update to the \texttt{q} lane
factor through distinct valuation exponents, so for \texttt{p != q} the order
of those two tracked updates should not change the projected state. The tracked
route proves that independent prime-lane updates commute componentwise. The
positive surfaces are the tracked G6 commuting theorem and the tracked above-threshold
coordinate-independence consumer. The old full \texttt{LaneOperator} interface
is not promoted; its universal
vanished-prime identity law blocks the nontrivial scaling route used by the
tracked theorem. Migration from the tracked route to the older universal
interface, or to downstream physics, remains uninhabited.

\subsection{E8/LILA: Local Completeness and Upstream Promotion Boundary}

The E8/LILA application uses the same spine in an enumeration setting. The
carrier is finite local root geometry: integer and sparse local shapes can be
enumerated, audited, and compared against the local semantic boundary. The
positive and boundary surfaces are the E8 integer/two-sparse completeness
theorem, the E8 upstream-promotion boundary, and the E8 upstream-promotion
audit. The positive result is local integer/two-sparse completeness under the
current audit boundary. The missing adapter is the upstream completion
constructor or promotion API, not a stronger physics claim by local enumeration
alone.

\begin{lemma}[Root-shell emergence as bounded local geometry]
Suppose a finite carrier region admits signed local coordinates, a norm or
shell predicate, and an admissibility test that retains only locally compatible
coordinate differences. Then the visible local states are organized first by
shell membership and adjacency, not by arbitrary list position. When the
admissible shell is finite, enumeration can certify completeness of that
bounded shell and detect duplicate or missing local representatives.
\end{lemma}

\begin{proof}
The norm/shell predicate partitions candidate local differences into finite
fibres. The compatibility test removes candidates outside the current local
geometry. Completeness is therefore a finite membership statement: every
admissible representative in the shell has been listed, and no listed
representative violates the predicate. This is the shape of the local
E8/LILA root-boundary surface. It does not supply an upstream E8 completion
token or a physical representation theorem.
\end{proof}

\section{Locality, Time, and Causality from Finite-Support Traversal}

\begin{definition}[Finite-support configuration]
Let \(I\) be a coordinate index set and let \(\mathcal A_i\) be the value
space at coordinate \(i\). A configuration is an admissible finite-support
assignment
\[
x\in\mathcal C\subseteq\prod_{i\in I}^{\mathrm{fin}}\mathcal A_i,
\qquad
\operatorname{supp}(x)=\{i\in I\mid x(i)\neq 0_i\}.
\]
\end{definition}

\begin{lemma}[Finite support gives a local boundary]
If \(x,y\in\mathcal C\) have finite support, then
\[
\Delta(x,y)=\{i\in I\mid x(i)\neq y(i)\}
\]
is finite.
\end{lemma}

\begin{proof}
\(\Delta(x,y)\subseteq
\operatorname{supp}(x)\cup\operatorname{supp}(y)\), and the right side is
finite. Thus any actual state difference has a finite coordinate witness.
\end{proof}

\begin{definition}[Local move and traversal history]
A move \(m:\mathcal C\rightharpoonup\mathcal C\) has finite footprint
\(\operatorname{foot}(m)\subset I\) when \(m(x)=y\) implies
\[
x|_{I\setminus\operatorname{foot}(m)}
=
y|_{I\setminus\operatorname{foot}(m)}.
\]
A traversal history is a sequence
\[
H=(x_0\xrightarrow{m_0}x_1\xrightarrow{m_1}x_2\cdots)
\]
with \(x_{k+1}=m_k(x_k)\). Its internal time is the occurrence order of this
sequence, not an assumed physical spacetime parameter.
\end{definition}

\begin{lemma}[Local moves preserve finite support]
If \(x\) has finite support and \(m\) has finite footprint, then \(m(x)\),
when defined, has finite support.
\end{lemma}

\begin{proof}
If \(y=m(x)\), then
\[
\operatorname{supp}(y)\subseteq
\operatorname{supp}(x)\cup\operatorname{foot}(m).
\]
The right side is finite.
\end{proof}

\begin{definition}[Traversal filtration and causal preorder]
Define
\[
S_n(H)=
\bigcup_{k\le n}\operatorname{supp}(x_k)
\cup
\bigcup_{k<n}\operatorname{foot}(m_k).
\]
Then \(S_0(H)\subseteq S_1(H)\subseteq\cdots\). An event
\(e_k=(k,m_k,x_k,x_{k+1})\) directly depends on \(e_j\) when \(j<k\) and the
later move inspects a coordinate affected or certified by the earlier move,
for example
\[
\operatorname{foot}(m_j)\cap\operatorname{insp}(m_k)\neq\varnothing.
\]
The reflexive-transitive closure gives an internal causal preorder subordinate
to traversal order.
\end{definition}

\begin{lemma}[Finite-speed propagation in traversal steps]
If the coordinate set carries a locally finite graph and every move footprint
lies in a bounded neighborhood of its inspection set, then after \(n\)
traversal steps the support of \(x_n\) lies in the \(n\)-step reachable region
from the initial support.
\end{lemma}

This is an internal locality theorem target. It does not yet assert
relativistic locality, light cones, or spacetime completion. The residual
obligations are the coordinate adjacency structure, admissible move class,
finite-branching law, dependence predicate, and history-equivalence criterion.

\begin{corollary}[Lorentz bridge as propagation-envelope target]
If local moves have a uniform propagation bound and admissible histories
preserve an invariant propagation envelope, then the reachability preorder
determines cone-shaped regions
\[
\mathcal C^+(e)=\{e'\mid e\preceq e'\},
\qquad
\mathcal C^-(e)=\{e'\mid e'\preceq e\}.
\]
A Lorentz-like structure is the additional target in which these cones are
represented by an invariant quadratic envelope, for example
\[
q(\Delta)=c^2\Delta t^2-\|\Delta x\|^2\ge 0,
\]
up to the accepted signature and calibration laws.
\end{corollary}

\begin{proof}
Bounded traversal gives finite reachable regions at each step. If the bound
is stable under admissible changes of local coordinates, then reachability
cannot be represented as an arbitrary graph property; it becomes an invariant
envelope of possible influence. The Lorentz target is the case where that
envelope is coordinatized by a nondegenerate quadratic cone. The present paper
does not inhabit that target; it names the adapter from finite traversal to
invariant cone geometry.
\end{proof}

\section{Gauge Transport from Local Sections}

\begin{definition}[Local section and transition]
Let \(C\) be the finite carrier and let \(\pi:P\to C\) be a bundle target. A
local chart over \(U_i\subseteq C\) is a section
\[
s_i:U_i\to P,
\qquad
\pi\circ s_i=\mathrm{id}_{U_i}.
\]
On overlaps \(U_{ij}=U_i\cap U_j\), a transition map is a gauge element
\[
g_{ij}:U_{ij}\to G,
\qquad
s_j(x)=s_i(x)\cdot g_{ij}(x),
\]
with target cocycle laws \(g_{ii}=1\), \(g_{ji}=g_{ij}^{-1}\), and
\(g_{ij}g_{jk}=g_{ik}\) on triple overlaps.
\end{definition}

The terminology in this section is the standard bundle, connection, and
curvature terminology used in gauge geometry
\cite{nakahara-2003,baez-muniain-1994,marsh-2016}; the DASHI contribution is
the carrier-side staging of which pieces are inhabited, targeted, or still
residual.

\begin{definition}[Connection and plaquette defect]
A local connection assigns an oriented edge \(e:x\to y\) a transport
\(U_e:F_x\to F_y\), or after trivialization an element \(U_e\in G\). For a
plaquette \(p\) with oriented boundary \(e_1+e_2-e_3-e_4\), define holonomy
\[
\operatorname{Hol}(p)=U_{e_4}^{-1}U_{e_3}^{-1}U_{e_2}U_{e_1},
\]
up to boundary convention, and discrete curvature defect
\[
\Omega(p)=\operatorname{Hol}(p)-1
\quad\text{or}\quad
F(p)\approx\log(\operatorname{Hol}(p)).
\]
\end{definition}

\begin{lemma}[Abelian transport reduces to a boundary sum]
When link transports commute, plaquette holonomy can be represented additively
by a cochain \(a:E_1\to A\):
\[
F(p)=d_1a(p)=a(e_1)+a(e_2)-a(e_3)-a(e_4),
\qquad
a\mapsto a+d_0\lambda.
\]
Curvature is gauge invariant when \(d_1d_0=0\).
\end{lemma}

\begin{definition}[Nonabelian target]
When transports do not commute, the order around the plaquette matters and
the smooth target becomes
\[
F=dA+A\wedge A,
\qquad
F_{\mu\nu}=
\partial_\mu A_\nu-\partial_\nu A_\mu+[A_\mu,A_\nu].
\]
This is a target adapter, not a Maxwell or Yang-Mills closure. The residual
obligations are carrier patches, local sections, transition maps, cocycle
laws, oriented cells, finite incidence, \(d_1d_0=0\), a nondegenerate curvature
carrier, and field-equation witnesses.
\end{definition}

\begin{corollary}[Why the abelian/nonabelian split has this target shape]
If the transport group is abelian, the plaquette defect is a
boundary-of-boundary object: \(F=dA\), gauge change is
\(A\mapsto A+d\lambda\), and the homogeneous law \(dF=0\) follows from
\(d^2=0\). This is the DASHI route to an EM-shaped abelian sector. If the
transport group is nonabelian, additive boundary calculus is not stable under
composition; the commutator correction is the minimal local term that records
the ordering defect, giving the Yang-Mills-shaped curvature
\(F_A=dA+A\wedge A\). The field equations \(d{*F}=J\) and
\(d_A{*F_A}=J\) are not obtained from this split alone; they require a
Hodge/source/action adapter.
\end{corollary}

\begin{theorem}[Plaquette ordering motivates the Yang-Mills curvature target]
Let \(U_\mu(x)=\exp(\epsilon A_\mu(x))\) and
\(U_\nu(x)=\exp(\epsilon A_\nu(x))\) be infinitesimal transports in two local
directions. The second-order plaquette holonomy has the form
\[
U_\nu(x)^{-1}U_\mu(x+\nu)^{-1}U_\nu(x+\mu)U_\mu(x)
=
I+\epsilon^2\bigl(
\partial_\mu A_\nu-\partial_\nu A_\mu+[A_\mu,A_\nu]
\bigr)+O(\epsilon^3),
\]
up to the chosen orientation convention. Thus the nonabelian curvature target
is not an optional embellishment: the commutator term is the first
order-sensitive local obstruction to plaquette closure.
\end{theorem}

\begin{proof}
Expand each edge transport to second order and compare the two orders around
the plaquette. The derivative terms record the change of the local connection
along neighboring edges. The only second-order term that survives solely
because the two transports fail to commute is \([A_\mu,A_\nu]\). In the
abelian case this term vanishes and the curvature reduces to \(dA\). The
residual obligations are the accepted smooth-limit adapter, representation,
action/source law, Hodge operator, and calibrated field equation.
\end{proof}

\section{QFT and Operator Surfaces from Filtered Quotients}

Let \(A_{\mathrm{fs}}\) denote the selected finite-support operator surface
over the current G3 wave-function carrier \(W\):
\[
W=\mathrm{SelectedG3State}\to\mathbb Q,
\qquad
A_{\mathrm{fs}}\subseteq \operatorname{End}(W).
\]
Let \(F_n\) be the selected filtration piece. The associated graded target is
not \(F_n\) itself but
\[
\operatorname{gr}_n(F)=F_n/F_{n-1}.
\]
We use standard operator-algebra notation here
\cite{haag-1996,reed-simon-1980,weinberg-1995}; the constructed object is the
quotient-descent condition, not a field algebra.

\begin{definition}[Quotient descent law]
The quotient projection
\(\pi_n:F_n\to\operatorname{gr}_n(F)\) is well formed only with an equivalence
\[
A\sim_n B \quad\Longleftrightarrow\quad A-B\in F_{n-1}
\]
and the kernel law
\[
\pi_n(A)=\pi_n(B)
\quad\Longleftrightarrow\quad
A-B\in F_{n-1}.
\]
\end{definition}

\begin{definition}[Descended bracket target]
The local bracket is
\[
[A,B]=A\circ B-B\circ A.
\]
The associated-graded bracket target is
\[
[\pi_m(A),\pi_n(B)]_{\mathrm{gr}}
=
\pi_{m+n}([A,B]).
\]
This definition is representative-independent only if
\[
A\sim_m A',\;B\sim_n B'
\Rightarrow
[A,B]-[A',B']\in F_{m+n-1}.
\]
\end{definition}

\begin{definition}[Amplitude and commutator target]
A physical operator claim requires a structure map
\[
\Phi:\operatorname{gr}(F)\to \operatorname{End}(\mathcal H)
\]
that preserves the descended bracket. A canonical commutator is therefore a
residual target
\[
R_{\mathrm{CCR}}(i,j)=
[\Phi(X_i),\Phi(P_j)]-i\hbar\delta_{ij}I,
\]
with promotion only after \(R_{\mathrm{CCR}}(i,j)=0\) is inhabited under an
accepted state space, scalar/amplitude structure, identity, calibration, and
quotient-descent receipt.
\end{definition}

\begin{lemma}[Why the commutator is the obstruction object]
For operators that descend through \(\pi_n\), the expression
\(A\circ B-B\circ A\) is the least local witness of order-dependence: it
vanishes exactly when the two updates commute at the chosen projection grade.
Thus the QFT-facing bracket does not enter as a decorative analogy. It is the
obstruction to exchanging filtered updates after quotienting. The remaining
physics step is the representation map \(\Phi\) that turns that obstruction
calculus into a calibrated operator algebra over \(\mathcal H\).
\end{lemma}

The residual obligations are the quotient carrier, kernel-equivalence law,
finite-support product and commutator closure, bracket descent, amplitude
space \(\mathcal H\), bracket preservation, and physical calibration. QFT,
local QFT, field algebra, Born rule, and interaction closure remain open.

\section{Curvature and GR from Transport Defect}

The GR target remains
\[
G_{\mu\nu}+\Lambda g_{\mu\nu}
=
\frac{8\pi G}{c^4}T_{\mu\nu}.
\]
In this paper the equation is a target: each term specifies a carrier-side
construction still to be supplied. The left side requires
metric, connection, curvature, Ricci contraction, scalar curvature, Einstein
tensor, cosmological term, and a contracted-Bianchi law. The right side
requires an authority-backed matter/stress-energy tensor, physical-unit
normalization, and covariant conservation.
The curvature and Bianchi notation follows standard GR conventions
\cite{wald-1984,carroll-2019,misner-thorne-wheeler-1973}; DASHI uses it here
to name the exact adapter targets.

\begin{definition}[Transport defect]
Let \(X\) be a discrete carrier with vertices \(x\), oriented edges \(e:x\to
y\), and local fibres \(F_x\). A local transport assigns \(U_e:F_x\to F_y\).
For a path \(\gamma=e_n\cdots e_1\), let
\(U_\gamma=U_{e_n}\cdots U_{e_1}\). For a plaquette \(p\) based at \(x\),
define
\[
H_p=U_{\partial p}:F_x\to F_x,
\qquad
\Omega_p=H_p-\mathrm{id}_{F_x}.
\]
A non-flat curvature target requires a carrier where some plaquette defect is
nonzero and transforms covariantly under frame change.
\end{definition}

\begin{definition}[Connection and curvature target]
A discrete connection can be treated as primitive transport \(U_{x,\mu}\), or
by an expansion
\[
U_{x,\mu}=I+\epsilon\Gamma_\mu(x)+O(\epsilon^2).
\]
The covariant difference target is
\[
\nabla_\mu\phi(x)=U_{x,\mu}^{-1}\phi(x+\mu)-\phi(x).
\]
For a GR route, this must supply \(\Gamma^\rho_{\mu\nu}\), \(g_{\mu\nu}\),
\(g^{\mu\nu}\), metric compatibility \(\nabla_\lambda g_{\mu\nu}=0\), and
torsion-free symmetry where the Levi-Civita route is intended.

Curvature is the failure of covariant transport to commute:
\[
[\nabla_\mu,\nabla_\nu]V^\rho
=
R^\rho{}_{\sigma\mu\nu}V^\sigma
-T^\lambda{}_{\mu\nu}\nabla_\lambda V^\rho.
\]
The torsion-free GR target sets the torsion term to zero.
\end{definition}

\begin{definition}[Bianchi, contraction, and source]
The Bianchi targets are
\[
R^\rho{}_{[\sigma\mu\nu]}=0,
\qquad
\nabla_{[\lambda}R^\rho{}_{|\sigma|\mu\nu]}=0.
\]
Once \(R^\rho{}_{\sigma\mu\nu}\) exists, the Ricci and scalar contractions are
\[
R_{\mu\nu}=R^\rho{}_{\mu\rho\nu},
\qquad
R=g^{\mu\nu}R_{\mu\nu},
\]
and the Einstein tensor is
\[
G_{\mu\nu}=R_{\mu\nu}-\frac12 Rg_{\mu\nu}.
\]
The sourced equation then requires the same carrier, index convention, metric,
connection, units, and physical regime for \(G_{\mu\nu}\), \(\Lambda
g_{\mu\nu}\), and \(T_{\mu\nu}\).
\end{definition}

\begin{lemma}[Why the Einstein tensor is the natural source-facing target]
In the standard GR target, once a Levi-Civita route and contracted-Bianchi
adapter are supplied, the contracted Bianchi identity gives
\[
\nabla^\mu G_{\mu\nu}=0.
\]
Any sourced curvature equation coupled to a conserved stress-energy target
must therefore put a divergence-free curvature object on the left-hand side.
The cosmological term is compatible because \(\nabla^\mu g_{\mu\nu}=0\). This
does not prove the DASHI GR adapter; it explains why the adapter must target
\[
G_{\mu\nu}+\Lambda g_{\mu\nu}=\kappa T_{\mu\nu}
\]
rather than an arbitrary curvature expression. The value
\(\kappa=8\pi G/c^4\) remains a weak-field/Newtonian calibration obligation.
\end{lemma}

The residual obligations are a non-flat connection, curvature covariance,
finite-radius Bianchi identity, metric/inverse/contraction laws, contracted
Bianchi law, physical stress-energy interface, unit constants \(G,c,\Lambda\),
and an inhabited sourced equation.

\section{Measurement, Interference, Entropy, and Sector Splitting}

Let \(X\) be a carrier state space and let
\[
\pi:X\to Q
\]
be a projection, coarse-graining, or quotient map into an observational state
space. Two states are observationally equivalent when
\[
x\sim_\pi x' \Longleftrightarrow \pi(x)=\pi(x').
\]
Measurement is modeled first as selecting or stabilizing a quotient
representative in \(Q\), rather than as direct access to \(x\).

\begin{definition}[Residual fibre]
The residual fibre of \(x\) is
\[
R_\pi(x)=\pi^{-1}(\pi(x)).
\]
Uncertainty is represented by the residual equivalence class left by the
projection, pending a probability/amplitude law.
\end{definition}

\begin{theorem}[Incompatible projections force residual growth]
Let \(\pi_1:X\to Q_1\) and \(\pi_2:X\to Q_2\) be two projections. If there is
no projection \(\pi_{12}:X\to Q_{12}\) and maps \(r_i:Q_{12}\to Q_i\) such
that
\[
\pi_i=r_i\circ\pi_{12}
\quad(i=1,2),
\]
while preserving the required distinctions of both \(Q_1\) and \(Q_2\), then
the two observations admit no joint quotient representative. Sequential
measurement must therefore leave an order-dependent residual
\[
\begin{aligned}
D_{12}(x)
&=
\pi_2\!\left(s_1(\pi_1(x))\right)
-
\pi_1\!\left(s_2(\pi_2(x))\right),
\end{aligned}
\]
where \(s_i\) is any chosen section or representative rule. Nonzero
\(D_{12}\) is the projection-defect route to uncertainty and interference
targets.
\end{theorem}

\begin{proof}
A joint quotient representative is precisely a common refinement through
which both projections factor. If no such refinement preserves the demanded
distinctions, choosing a representative for one projection collapses
distinctions needed by the other. The lost information moves into a residual
fibre, and different section choices can produce different downstream
projected states. A probability or amplitude law is still required before
this becomes quantum measurement theory.
\end{proof}

\begin{definition}[Interference and entropy targets]
If unresolved branches \(x_i\) combine before projection,
\[
\Psi=\sum_i a_i x_i,
\qquad
\pi(\Psi)
\]
can depend on branch combination before quotient separation. Entropy may be
read projection-relatively as residual volume:
\[
S_\pi(q)\sim\log\mu(\pi^{-1}(q)),
\qquad
S_\pi=-\sum_i p_i\log p_i
\]
when branch weights \(p_i\) are supplied by the chosen formal model.
\end{definition}

\begin{definition}[Sector splitting]
Let \(\mathcal A\) be an algebra of transformations on \(X\). The
projection-compatible subalgebra is
\[
\mathcal A_\pi=\{a\in\mathcal A\mid \pi\circ a
\text{ is well defined on }Q\}.
\]
Commuting projected operations \([a,b]_\pi=0\) form abelian sector targets.
Noncommuting residual operations \([a,b]_\pi\neq0\) point toward nonabelian
targets. If a symmetry group \(G\) acts on \(X\) but projection preserves only
\(H\subseteq G\), then residual cosets \(G/H\) label directions no longer
treated equivalently by the projected description. This is a target mechanism
for sector splitting, not a derivation of the observed weak/strong/EM
decomposition.
\end{definition}

The residual obligations are the base category, the precise nature of \(\pi\),
branch-weight law, interference theorem, residual commutator, concrete group
\(G\), subgroup \(H\), representation data, and empirical calibration.

\section{Empirical Contact}

The primary empirical-contact surface used in this paper is the Gate 5
penguin-decay \(P_5'\) lane. The canonical diagnostic records
\texttt{empiricalContactReached = true}: the selected \(b\to s\ell\ell\)
angular-observable lane binds a checksum-backed value-table route, a frozen
Standard Model baseline request, and a residual-comparison law. This is
empirical contact, not empirical promotion.

The bound observable is \(P_5'\) for \(B^0\to K^{*0}\mu^+\mu^-\). The reachable
machine-readable LHCb route is HEPData
\texttt{hepdata.167733.v1/t2}, locally bound to
\texttt{scripts/data/hepdata/ins3094698/table-config-2-results-record-data-167733-1911073.json}
with SHA256
\texttt{8ee74f4e774889eced2090fe60bbdaf681dd327dc1a349add992e038c7f62623}.
The surrounding digest-bound artifacts include the total, statistical, and
systematic correlation JSON resources, the CDS supplementary archive, and the
2020 value-table/projection artifacts. The Standard Model baseline is the
\texttt{flavio 2.7.0} \(P_5'\) baseline with the Bharucha--Straub--Zwicky
(BSZ) form-factor source; its manifest is
\texttt{scripts/data/hepdata/penguin\_gate5\_p5prime\_sm\_baseline\_flavio\_2\_7\_0.json}.

The residual vector is classified as
\texttt{p5PrimeBorderlineAnomalyCandidate}. Using
\[
r=(\mathrm{obs}-\mathrm{SM})/
\sqrt{\mathrm{stat}^2+\mathrm{syst}^2+\mathrm{th}^2},
\]
the selected \(q^2\) bins record signed pulls
\(-2.8\sigma\) for \([4,6]\,\mathrm{GeV}^2/c^4\) and
\(-3.0\sigma\) for \([6,8]\,\mathrm{GeV}^2/c^4\). These two pulls make the lane
a borderline anomaly-candidate surface, but the accepted residual predicate is
not inhabited:
\texttt{acceptedResidualCandidate = false}. The full \(P_5'\) covariance
matrix required for covariance-whitened promotion remains missing, and the
accepted authority, complete freeze, data, and controlled-theory prerequisites
are not all present.

\subsection*{Typed Residual Basis Decomposition}

The Drell-Yan material is retained as a methodology demonstration rather than
the primary empirical result. Typed Residual Basis Decomposition (TRBD)
records
\[
\begin{aligned}
\text{residual}
&\longrightarrow
\text{projection onto a typed basis}
\longrightarrow
\text{named obstruction}
\\
&\longrightarrow
\text{coverage and perpendicular residual}
\longrightarrow
\text{promotion decision}.
\end{aligned}
\]
The reusable core is \texttt{DASHI.Core.TypedResidualBasisDecomposition}. The
Drell-Yan strict-log diagnostic instantiates it with the structural basis
\(\{1,\log(\varphi^*)\}\). For \texttt{sigmaDashiV4Predictor}, the raw
strict-log \texttt{chi2/dof} is \texttt{3180.211733150705}, the perpendicular
\texttt{chi2/dof} is \texttt{111.96455543013676}, and the basis coverage is
\texttt{0.968705212853035}. The coverage is high, but the perpendicular
residual still fails the strict threshold \texttt{2.0}, so the TRBD status is
\texttt{obstructionTypedPartial}; no promotable TRBD receipt is constructed.

Operationally, the Drell-Yan TRBD diagnostic decomposes
\(\Delta\sigma_i=\sigma_i^{\rm data}-\sigma_i^{\rm theory}\) at the transition
regions around \(\varphi^*\approx0.14\) and \(\varphi^*\approx0.84\), using
log-linear, log-quadratic, and log-cubic projections. The observed sign change
near \(\varphi^*=0.1395\) is not treated as endpoint antisymmetry; it is
recorded as evidence for a two-transition residual basis. The implementation
surface is \texttt{dy\_slope\_decomposition.py}, while
\texttt{DrellYanLogLinearShapeLawReceipt.agda} binds the strict-log
obstruction and keeps the stronger lane blocked on a native Z+jet,
event-level \(\varphi^*\) provider.

The older bounded W3 below-\(Z\) Drell-Yan comparison-law receipt remains a
non-primary example of bounded contact accounting: over CMS-SMP-20-003 /
HEPData \texttt{ins2079374/t43}, using the \texttt{ins2079374/t44} covariance
convention, it records \texttt{chi2/dof = 2.1565191176}, threshold
\texttt{4.0}, and mean prediction/data \texttt{0.9941233097}. This is a TRBD
and comparison-law demonstration, not the paper's primary empirical result,
and it has no above-\(Z\), strict \texttt{chi2/dof <= 2}, or accepted-authority
promotion.

The shared boundary is fail-closed: \(P_5'\) reaches empirical contact, while
Drell-Yan illustrates typed residual decomposition. Neither lane constructs an
accepted external empirical authority token or licenses discovery, new
physics, full empirical adequacy, W4/W5 convention closure, or physical
calibration.

\section{Compression, Semantic Geometry, and Cross-Domain Spine}

UFT addresses compress repeated refinement structure. Two states may share a
motif at grade \(\Omega\) only when the distinctions required below \(\Omega\)
remain recoverable. A coarse address names a region of structure; descending
the trie adds precision, distinguishes branches, and exposes incompatibilities.
Branch distance is ultrametric, branch membership is projection-grade
dependent, and unresolved branches can remain live without binary collapse.

\begin{theorem}[Compression admissibility at a projection grade]
Let \(\Xi\) be a semantic state, let \(\Omega\) be the active projection grade,
and let \(\Pi_{\mathrm{compress}}(\Xi)\) be a compressed representative. The
compression is admissible at grade \(\Omega\) only if every finer grade
\(\Omega'\leq\Omega\) has an explicit recoverability witness:
\[
\forall \Omega'\leq\Omega,\quad
\mathrm{RequiredDistinctions}(\Omega')
\subseteq
\mathrm{Recoverable}_{\Omega'}
\bigl(\Pi_{\mathrm{compress}}(\Xi)\bigr).
\]
\end{theorem}

The compression admissibility receipt inhabits this condition at motif grade
for exact, PNF, and motif recoverability witnesses. It also records that the
residual carrier is preserved, CID identity and semantic identity are
separated, and false semantic collapse is excluded.

The same receipt also records the MDL form of admissibility:
\[
L_{\rm MDL}(T(s))\leq L_{\rm MDL}(s).
\]
Thus admissible carrier transport is description-length non-increasing at the
chosen projection grade. Cone monotonicity and max-severity joins are
downstream of this condition. Vacuum energy and cosmological-constant terms
remain MDL-calibration open obligations: no receipt currently supplies a
description-length bound that would promote those physical claims. The receipt
status is diagnostic: it does not construct an external ITIR runtime authority
token and does not identify byte-exact CID equality with semantic equivalence.

UFT-C, lattice-logic, null/hinge codes, SWAR implementation machinery,
Base369/trit notation, Hensel-style refinement, and Markov-after-quotient
readings are routed as background or appendix material unless a repo-local
surface promotes a stronger claim.

The cross-domain variational-spine boundary records a common typed schema
\((X,\delta,\pi,\mathrm{defect},\mathrm{gate},\mathrm{observation},
\mathrm{symmetry})\) shared by several non-promoting target rows. Paper 1 does
not infer chemistry closure, biological prediction, perceptual empirical fit,
or universal closure from that boundary.

\section{Closure Frontier}

The closure frontier records the exact residual obligations remaining for
stronger physical claims.

\begin{longtable}{p{0.16\linewidth}p{0.29\linewidth}p{0.27\linewidth}p{0.17\linewidth}}
\toprule
Gate & Current surface & First missing primitive/token & Unsupported stronger claim \\
\midrule
G2 & Direction-indexed SFGC schema adoption plus standalone prime-lattice \(\mathrm{d}^2=0\) targets & Prime-lattice commutativity/boundary-of-boundary law, SFGC action bridge, curvature/field-equation carrier, and promotion receipts & Maxwell closure, gauge-field recovery, physical curvature \\
G3 & Full depth-9 curvature and finite \(D \star F = 0\) inhabited & Continuum analytic promotion remains external & Continuum Yang--Mills / PDE authority, arbitrary analytic field-equation closure \\
G4 & Finite sourced Einstein interface inhabited with Wald authority boundary recorded & Continuum Levi-Civita uniqueness remains external & Continuum Einstein equations, Schwarzschild and GR recovery, unqualified sourced GR \\
G5 & Empirical contact reached with \texttt{empiricalContactReached = true}; \(P_5'\) \(-2.8/-3.0\) sigma candidate for \texttt{hepdata.167733.v1/t2}; accepted residual candidate is false & Full covariance is missing; no accepted residual authority token & Accepted B-anomaly authority, empirical closure, Standard Model deviation claim \\
G6 & Finite endomorphism, naturality, hexagon, statistics, and conjugate-zigzag receipts inhabited & Arbitrary DHR symmetric monoidal structure plus DR authority/compact reconstruction & Full DHR category, compact gauge group reconstruction, Standard Model sector promotion \\
G7 & Identity and eigenbasis transport receipts inhabited & Physical Yukawa construction from DHR sector remains blocked & Physical Yukawa matrices, CKM/PMNS derivation, fermion-mass closure \\
NS & Finite Navier--Stokes ladder inhabited & Continuum BKM/enstrophy lift remains blocked & Clay Navier--Stokes regularity / global smoothness claim \\
Information paradox & Cross-gate obstruction surface recorded & Joint accepted GR/QFT/thermodynamic authority and adapter alignment & Resolved black-hole information paradox claim \\
E8/LILA & Local integer/root boundary & Upstream completion constructor/API & Upstream E8 completion \\
W3 & Bounded Drell-Yan receipt & Accepted W3 evidence-authority token & Accepted empirical authority \\
W4/W5 & Request/diagnostic surfaces & Accepted DY convention and correction receipts & Matter field, stress-energy, sourced physics \\
DCHoTT bridge & \texttt{LeviCivitaBridge.agda} plus B0.1/B0.2/B0.3 sockets; DASHI-side B0.1 pro-object semantics & DCHoTT formal-D-space adapter; real-valued ultrametric completion; refinement-stable frame/metric tower; G-structure reduction; torsion-free/Levi-Civita specialization & Levi-Civita uniqueness, GR promotion \\
AQFT net & AQFT/Klein-Gordon receipts plus carrier-algebra quotient sockets & Constructive local algebra quotient, filtered colimit, Cauchy time-slice receipt, GNS/vacuum adapter, Born-rule derivation & Interacting QFT, Standard Model \\
GRQFT / terminal boundary & Weak terminal grammar modulo explicit postulates; terminal promotion blocked; DASHI-internal ST3/Banach mass-gap closure and conditional Moonshine \texttt{laneDimension} bridge recorded separately & B0 plus AQFT net, DHR/SM comparison, Clay/external mass-gap authority, \texttt{DASHIPrimeLaneEquiv}, final adapter acceptance, stress-energy/Wald authority & Unqualified terminal GRQFT/TOE claim \\
GR & Future adapter path & Non-flat connection, curvature, Ricci, sourced interface & Einstein tensor, Schwarzschild, continuum GR \\
Cross-domain spine & Shared projection-defect / MDL schema & Calibration, universality, domain receipts & Cross-domain recovery \\
Provenance & JMD/SSP/archive/source routing & Repo-local receipts or accepted authority & Provenance-as-proof \\
\bottomrule
\end{longtable}

Every row identifies a residual defect: the exact bridge, authority, quotient
law, adapter, or calibration object required for a stronger projection.


\section{Reproducibility}

The canonical manuscript, claim ledger, source ledger, archive-topic plan, and
relative archive-location reference are included in the accompanying
repository. Exact repository paths are listed in the supplemental repository
materials rather than in the main prose.

The narrow verification surface for this draft is:
\begin{verbatim}
plantuml -checkonly Docs/*.puml
./scripts/render_docs_diagrams.sh
git diff --check -- Docs/papers/drafts/PaperDraftWorkingFolder Docs README.md
\end{verbatim}

\subsection*{Receipt Index}

Agda identifiers are kept out of explanatory prose except where the exact name
is the object under discussion. The table below gives the exact receipt names
for the human-readable labels used in the body.

\begin{longtable}{p{0.20\linewidth}p{0.34\linewidth}p{0.34\linewidth}}
\toprule
Receipt label & Module & Agda identifier \\
\midrule
G2 schema extension &
\scriptsize\ttfamily\seqsplit{DASHI.Physics.Closure.G2SFGCGaugeFieldSchemaExtension} &
\scriptsize\ttfamily\seqsplit{canonicalG2SFGCGaugeFieldSchemaExtensionReceipt} \\
G2 no-missing witness &
\scriptsize\ttfamily\seqsplit{DASHI.Physics.Closure.G2SFGCGaugeFieldSchemaExtension} &
\scriptsize\ttfamily\seqsplit{noRemainingG2SFGCSchemaExtensionMissing} \\
G3 selected subtraction support &
\scriptsize\ttfamily\seqsplit{DASHI.Physics.Closure.G3P2LimitConvergenceSurface} &
\scriptsize\ttfamily\seqsplit{selectedFiniteSupportOperatorSubtractionCertifiedSupport} \\
G3 selected subtraction witness &
\scriptsize\ttfamily\seqsplit{DASHI.Physics.Closure.G3P2LimitConvergenceSurface} &
\scriptsize\ttfamily\seqsplit{selectedFiniteSupportOperatorSubtractionSupportWitnessFromCertified} \\
G3 certified-support record &
\scriptsize\ttfamily\seqsplit{DASHI.Physics.Closure.G3P2LimitConvergenceSurface} &
\scriptsize\ttfamily\seqsplit{G3SelectedSubtractionCertifiedSupport} \\
G6 tracked commuting theorem &
\scriptsize\ttfamily\seqsplit{DASHI.Physics.Closure.G6CrossLaneCommutingTheorem} &
\scriptsize\ttfamily\seqsplit{canonicalG6OfficialTrackedCrossLaneCommutingTheorem} \\
G6 above-threshold independence &
\scriptsize\ttfamily\seqsplit{DASHI.Physics.Closure.G6AboveThresholdIndependence} &
\scriptsize\ttfamily\seqsplit{canonicalG6OfficialTrackedAboveThresholdCoordinateIndependence} \\
E8 integer/two-sparse completeness &
\scriptsize\ttfamily\seqsplit{DASHI.Algebra.Trit.E8RootEnumeration} &
\scriptsize\ttfamily\seqsplit{integerIndexedRootsCompleteForTwoSparseShapeTheorem} \\
E8 upstream-promotion boundary &
\scriptsize\ttfamily\seqsplit{DASHI.Algebra.Trit.E8RootEnumeration} &
\scriptsize\ttfamily\seqsplit{canonicalE8RootEnumerationCompletePromotionBoundary} \\
E8 upstream-promotion audit &
\scriptsize\ttfamily\seqsplit{DASHI.Physics.Closure.LilaE8RootEnumerationNoDuplicatesSurface} &
\scriptsize\ttfamily\seqsplit{canonicalE8UpstreamCompleteReceiptPromotionAudit} \\
W3 bounded comparison law &
\scriptsize\ttfamily\seqsplit{DASHI.Physics.Closure.HEPDataW3ComparisonLawReceipt} &
\scriptsize\ttfamily\seqsplit{canonicalHEPDataW3ComparisonLawReceipt} \\
Drell-Yan adjacent-ratio lane &
\scriptsize\ttfamily\seqsplit{DASHI.Physics.Closure.DrellYanAdjacentRatioEmpiricalLaneReceipt} &
\scriptsize\ttfamily\seqsplit{canonicalDrellYanAdjacentRatioEmpiricalLaneReceipt} \\
Drell-Yan strict-log diagnostic &
\scriptsize\ttfamily\seqsplit{DASHI.Physics.Closure.DrellYanStrictLogLinearSubspaceReceipt} &
\scriptsize\ttfamily\seqsplit{canonicalDrellYanStrictLogLinearSubspaceReceipt} \\
Compression admissibility &
\scriptsize\ttfamily\seqsplit{DASHI.Physics.Closure.CompressionAdmissibilityReceipt} &
\scriptsize\ttfamily\seqsplit{canonicalCompressionAdmissibilityReceipt} \\
DCHoTT import shim &
\scriptsize\ttfamily\seqsplit{DASHI.Geometry.DCHoTTImportShim} &
\scriptsize\ttfamily\seqsplit{canonicalDCHoTTImportShimReceipt} \\
DCHoTT Levi-Civita bridge &
\scriptsize\ttfamily\seqsplit{DASHI.Geometry.LeviCivitaBridge} &
\scriptsize\ttfamily\seqsplit{canonicalLeviCivitaBridgeReceipt} \\
DCHoTT B0 obligation index &
\scriptsize\ttfamily\seqsplit{DASHI.Geometry.DCHoTTBridgeObligationIndex} &
\scriptsize\ttfamily\seqsplit{canonicalDCHoTTBridgeObligationIndex} \\
Supersingular-prime lane bridge &
\scriptsize\ttfamily\seqsplit{DASHI.Physics.Moonshine.SupersingularPrimeLaneBridge} &
\scriptsize\ttfamily\seqsplit{canonicalSupersingularPrimeLaneBridgeReceipt} \\
B0.1 pro-object semantics &
\scriptsize\ttfamily\seqsplit{DASHI.Geometry.ProObjectSemantics} &
\scriptsize\ttfamily\seqsplit{canonicalB01ProObjectSemanticsReceipt} \\
AQFT carrier algebra quotient surface &
\scriptsize\ttfamily\seqsplit{DASHI.Physics.QFT.AQFTCarrierAlgebraQuotientSurface} &
\scriptsize\ttfamily\seqsplit{canonicalAQFTCarrierAlgebraQuotientSurface / canonicalDepthFilteredLocalAlgebraSurface / canonicalCauchyEvolutionReceiptTarget} \\
DHR gauge receipt surface &
\scriptsize\ttfamily\seqsplit{DASHI.Physics.QFT.DHRGaugeReceiptSurface} &
\scriptsize\ttfamily\seqsplit{canonicalDHRGaugeReceiptSurface} \\
DHR finite-prime local endomorphism category ledger &
\scriptsize\ttfamily\seqsplit{DASHI.Physics.QFT.Gate6FinitePrimeLaneLocalEndomorphismCategory} &
\scriptsize\ttfamily\seqsplit{canonicalGate6FinitePrimeLaneLocalEndomorphismCategoryLedger} \\
DHR tensor-dual reconstruction boundary &
\scriptsize\ttfamily\seqsplit{DASHI.Physics.QFT.DHRTensorDualGroupReconstruction} &
\scriptsize\ttfamily\seqsplit{canonicalDHRTensorDualGroupReconstructionReceipt} \\
DHR finite-prime statistics/braiding progress &
\scriptsize\ttfamily\seqsplit{DASHI.Physics.QFT.DHRTensorDualGroupReconstruction} &
\scriptsize\ttfamily\seqsplit{canonicalFinitePrimeLaneTensorBraidingProgressReceipt} \\
DHR hexagon fail-closed coherence &
\scriptsize\ttfamily\seqsplit{DASHI.Physics.QFT.DHRHexagonObligation} &
\scriptsize\ttfamily\seqsplit{canonicalDHRHexagonFailClosedCoherenceReceipt} \\
DHR finite-prime hexagon target receipts &
\scriptsize\ttfamily\seqsplit{DASHI.Physics.QFT.DHRHexagonObligation} &
\scriptsize\ttfamily\seqsplit{canonicalDHRFinitePrimeLaneLeftRightHexagonTargetReceipts} \\
DHR finite-prime conjugate dual receipt &
\scriptsize\ttfamily\seqsplit{DASHI.Physics.QFT.FinitePrimeLaneConjugateDualReceipts} &
\scriptsize\ttfamily\seqsplit{canonicalFinitePrimeLaneConjugateDualReceipt} \\
DHR--SM compatibility ledger boundary &
\scriptsize\ttfamily\seqsplit{DASHI.Physics.QFT.FinitePrimeLaneDHRSMCompatibilityLedger} &
\scriptsize\ttfamily\seqsplit{canonicalFinitePrimeLaneDHRSMCompatibilityLedger} \\
Adapter irreducibility no-go index &
\scriptsize\ttfamily\seqsplit{DASHI.Physics.Closure.AdapterIrreducibilityNoGoIndex} &
\scriptsize\ttfamily\seqsplit{canonicalAdapterIrreducibilityNoGoIndex} \\
AQFT net receipt &
\scriptsize\ttfamily\seqsplit{DASHI.Physics.Closure.AQFTNetReceipt} &
\scriptsize\ttfamily\seqsplit{canonicalAQFTNetReceipt} \\
Klein-Gordon AQFT witness &
\scriptsize\ttfamily\seqsplit{DASHI.Physics.Closure.KleinGordonAQFTReceipt} &
\scriptsize\ttfamily\seqsplit{canonicalKleinGordonAQFTReceipt} \\
Interacting QFT boundary &
\scriptsize\ttfamily\seqsplit{DASHI.Physics.Closure.InteractingQFTBoundaryReceipt} &
\scriptsize\ttfamily\seqsplit{canonicalInteractingQFTBoundaryReceipt} \\
Terminal open-problem status &
\scriptsize\ttfamily\seqsplit{DASHI.Physics.Closure.TerminalOpenProblemStatusSurface} &
\scriptsize\ttfamily\seqsplit{canonicalTerminalOpenProblemStatusSurface} \\
GRQFT terminal composition boundary &
\scriptsize\ttfamily\seqsplit{DASHI.Physics.Closure.GRQFTTerminalCompositionBoundary} &
\scriptsize\ttfamily\seqsplit{canonicalGRQFTCompositionBoundary} \\
Navier-Stokes weak solution interface receipt &
\scriptsize\ttfamily\seqsplit{DASHI.Physics.Closure.NavierStokesWeakSolutionInterface} &
\scriptsize\ttfamily\seqsplit{canonicalNavierStokesWeakSolutionInterfaceReceipt} \\
Navier-Stokes regularity tower boundary receipt &
\scriptsize\ttfamily\seqsplit{DASHI.Physics.Closure.NavierStokesRegularityTowerReceipt} &
\scriptsize\ttfamily\seqsplit{canonicalNavierStokesRegularityTowerReceipt} \\
Wald GR authority fail-closed receipt &
\scriptsize\ttfamily\seqsplit{DASHI.Physics.Closure.WaldGRAuthorityReceipt} &
\scriptsize\ttfamily\seqsplit{canonicalWaldGRAuthorityReceipt} \\
P5-prime LHCb checksum/residual surface &
\scriptsize\ttfamily\seqsplit{DASHI.Physics.Closure.PenguinDecayLHCbChecksumAcceptedResidualReceipt} &
\scriptsize\ttfamily\seqsplit{canonicalLHCbChecksumAcceptedResidualSurfaceReceipt} \\
P5-prime residual comparison law boundary &
\scriptsize\ttfamily\seqsplit{DASHI.Physics.Closure.PenguinDecayResidualComparisonLaw} &
\scriptsize\ttfamily\seqsplit{canonicalPenguinDecayResidualComparisonLaw} \\
P5-prime SM baseline authority request diagnostic &
\scriptsize\ttfamily\seqsplit{DASHI.Physics.Closure.PenguinDecaySMBaselineAuthority} &
\scriptsize\ttfamily\seqsplit{canonicalSMBaselineAuthorityRequestDiagnostic} \\
P5-prime Wilson coefficient authority boundary &
\scriptsize\ttfamily\seqsplit{DASHI.Physics.Closure.PenguinDecayWilsonCoefficientAuthority} &
\scriptsize\ttfamily\seqsplit{canonicalPenguinDecayWilsonCoefficientAuthority} \\
P5-prime empirical candidate diagnostic &
\scriptsize\ttfamily\seqsplit{DASHI.Physics.Closure.PenguinDecayEmpiricalCandidateDiagnostic} &
\scriptsize\ttfamily\seqsplit{canonicalPenguinDecayEmpiricalCandidateDiagnostic} \\
P5-prime C9/C10 prediction target boundary &
\scriptsize\ttfamily\seqsplit{DASHI.Physics.Closure.PenguinDecayC9C10P5PrimePredictionTargetReceipt} &
\scriptsize\ttfamily\seqsplit{canonicalC9C10P5PrimePredictionTargetReceipt} \\
P5-prime carrier-derived C9 constraint target &
\scriptsize\ttfamily\seqsplit{DASHI.Physics.Closure.PenguinDecayCarrierDerivedC9ConstraintTargetReceipt} &
\scriptsize\ttfamily\seqsplit{canonicalPenguinDecayCarrierDerivedC9ConstraintTargetReceipt} \\
CKM prediction frontier receipt &
\scriptsize\ttfamily\seqsplit{DASHI.Physics.Closure.CKMPredictionFrontierReceipt} &
\scriptsize\ttfamily\seqsplit{canonicalCKMPredictionFrontierReceipt} \\
Yukawa carrier receipt boundary &
\scriptsize\ttfamily\seqsplit{DASHI.Physics.Closure.YukawaFromCarrier} &
\scriptsize\ttfamily\seqsplit{canonicalYukawaFromCarrier} \\
Carrier Yukawa ratio target receipt &
\scriptsize\ttfamily\seqsplit{DASHI.Physics.Closure.CarrierYukawaRatioTargetReceipt} &
\scriptsize\ttfamily\seqsplit{canonicalCarrierYukawaRatioTargetReceipt} \\
Information paradox cross-gate obstruction &
\scriptsize\ttfamily\seqsplit{DASHI.Physics.Closure.InformationParadoxCrossGateObstruction} &
\scriptsize\ttfamily\seqsplit{canonicalInformationParadoxCrossGateObstruction} \\
Quantum-gravity prediction frontier receipt &
\scriptsize\ttfamily\seqsplit{DASHI.Physics.Closure.InformationParadoxCrossGateObstruction} &
\scriptsize\ttfamily\seqsplit{canonicalBeyondCurrentRepoQuantumGravityPredictionFrontierReceipt} \\
YM depth-9 connection and curvature receipt &
\scriptsize\ttfamily\seqsplit{DASHI.Physics.Closure.Depth9ConnectionAndCurvature} &
\scriptsize\ttfamily\seqsplit{canonicalDepth9ConnectionAndCurvatureReceipt} \\
YM finite nonzero curvature computation &
\scriptsize\ttfamily\seqsplit{DASHI.Physics.Closure.Depth9ConnectionAndCurvature} &
\scriptsize\ttfamily\seqsplit{canonicalChainB1ExplicitNonzeroYMCurvatureFiniteComputation} \\
YM Gate 3 discrete-geometry receipt &
\scriptsize\ttfamily\seqsplit{DASHI.Physics.Closure.YangMillsGate3DiscreteGeometryReceipt} &
\scriptsize\ttfamily\seqsplit{canonicalGate3DiscreteGeometryReceipt} \\
YM field-equation obstruction boundary &
\scriptsize\ttfamily\seqsplit{DASHI.Physics.Closure.YangMillsFieldEquationObstruction} &
\scriptsize\ttfamily\seqsplit{canonicalYangMillsFieldEquationObstruction} \\
YM field-equation receipt boundary &
\scriptsize\ttfamily\seqsplit{DASHI.Physics.Closure.YangMillsFieldEquationReceipt} &
\scriptsize\ttfamily\seqsplit{canonicalYangMillsFieldEquationReceipt} \\
External formal import roadmap &
\scriptsize\ttfamily\seqsplit{DASHI.Physics.Closure.ExternalFormalImportRoadmapReceipt} &
\scriptsize\ttfamily\seqsplit{canonicalExternalFormalImportRoadmapReceipt} \\
Cross-domain variational spine &
\scriptsize\ttfamily\seqsplit{DASHI.Physics.Closure.CrossDomainVariationalSpine} &
\scriptsize\ttfamily\seqsplit{canonicalCrossDomainVariationalSpineBoundary} \\
\bottomrule
\end{longtable}

\section{What This Paper Does Not Claim}

This paper presents a constructive carrier derivation architecture and several
partial inhabited surfaces. It does not claim completed empirical closure,
sourced non-flat GR, GRQFT unification, upstream E8 completion, accepted W3/W4/W5
authority, or completed physics unification. Stronger readings remain
conditioned on the residual obligations listed in the closure frontier. In
particular, weak terminal status modulo named postulates is not an unqualified
terminal theorem.

\section*{Appendix A: Glossary and Naming Conventions}
\renewcommand{\arraystretch}{1.12}
\begin{longtable}{p{0.24\linewidth}p{0.66\linewidth}}
\toprule
Term & Meaning in this paper \\
\midrule
G2 lane & DASHI-native direction-indexed carrier schema for local
plaquette/traversal constructions. It is not identified with the exceptional
Lie group \(G_2\) unless such an equivalence is explicitly inhabited. \\
G3 lane & DASHI-native filtered finite-support operator lane. The current
positive result is selected finite-support subtraction support; associated
graded quotient descent remains open. \\
G6 lane & DASHI-native tracked cross-lane commuting lane over
\texttt{FactorVec} valuation coordinates. It is not a universal
\texttt{LaneOperator} theorem. \\
E8/LILA & Local finite root-shell enumeration and audit surfaces associated
with bounded E8-style root geometry. Upstream E8 completion is separate. \\
Surface & A projection context with a visible carrier, admissibility gate,
projection map, and residual defect. \\
Inhabited & Constructed or proved in the current typed surface. \\
Target & A specified mathematical obligation shape that is not yet inhabited. \\
Receipt & A mechanized or repo-local witness record for a bounded claim. \\
Promotion & Admissible elevation of a claim to the status allowed by its
evidence surface. \\
Residual / defect & Residual structure after projection. \\
Obstruction & A named missing bridge, law, token, adapter, or authority value
that blocks a stronger claim. \\
Calibration & The units, constants, empirical convention, or source-coupling
normalization needed for a physical claim. \\
Projection-defect decomposition & The reading \(x=P_s(x)+D_s(x)\), where
\(P_s(x)\) is the admitted projection and \(D_s(x)\) is residual structure. \\
UFT & Ultrametric Floating Trie: a hierarchical refinement-address structure
where shared prefixes induce ultrametric distance. \\
\texttt{FactorVec} & Prime-lane valuation vector indexed by the tracked prime
coordinate set. \\
\bottomrule
\end{longtable}
\renewcommand{\arraystretch}{1.0}

\section*{Appendix B: Symbol and Object Index}
\renewcommand{\arraystretch}{1.12}
\begin{longtable}{p{0.24\linewidth}p{0.66\linewidth}}
\toprule
Symbol & Meaning \\
\midrule
\(\mathcal X\) & DASHI carrier state space. \\
\(x\) & Carrier state. \\
\((a,\nu,\tau,d,\ell,r)\) & Address, valuation vector, ternary coordinate,
refinement depth, active lane, and receipt/projection status. \\
\(\tau\in\{-1,0,+1\}\) & Primitive ternary coordinate: opposed, unresolved,
affirmed. \\
\(P_s\) & Projection admitted by surface \(s\). \\
\(D_s(x)\) & Residual defect of \(x\) at surface \(s\). \\
\(\rho(x,y)\) & Ultrametric distance between carrier states. \\
\(F_n\) & Filtration layer. \\
\(\operatorname{gr}_n(F)\) & Associated graded target \(F_n/F_{n-1}\). \\
\(U_e\) & Transport operator on an oriented edge. \\
\(\operatorname{Hol}(p)\) & Plaquette holonomy. \\
\(\Omega(p)\) & Plaquette transport defect. \\
\(A,F_A\) & Connection one-form and nonabelian curvature target. \\
\(\nabla\) & Connection/covariant derivative target. \\
\(R^\rho{}_{\sigma\mu\nu}\) & Curvature target. \\
\(R_{\mu\nu},R\) & Ricci and scalar curvature targets. \\
\(G_{\mu\nu}\) & Einstein tensor target. \\
\(T_{\mu\nu}\) & Stress-energy/source target. \\
\(\pi:X\to Q\) & Observation or quotient projection. \\
\(R_\pi(x)\) & Residual fibre \(\pi^{-1}(\pi(x))\). \\
\(H\subseteq G\) & Preserved subgroup under sector projection. \\
\(G/H\) & Residual coset sector. \\
\bottomrule
\end{longtable}
\renewcommand{\arraystretch}{1.0}

\section*{Appendix C: Origins from Ternary Traversal to Carrier Geometry}
\setcounter{figure}{0}
\renewcommand{\thefigure}{A.\arabic{figure}}
\renewcommand{\theHfigure}{appendix.A.\arabic{figure}}

The earliest DASHI experiments used cyclic ternary and nonary carriers. A
trit distinguished positive, negative, and unresolved states; larger wrapped
state spaces represented transition, tension, overflow, and refinement without
collapsing the unresolved case into a binary decision. Later dialectical
language in the archive corresponds to the same mechanism: a transition may
carry unresolved residue, and the current projection records that residue as
\(D_s(x)\).

The geometric intuition then moves from a one-dimensional cycle to traversal
through a cube-like state space. Rubik-like cube, hypercube, and voxel language
is useful only as intuition for finite-coordinate movement: each coordinate
records a local trit, phase, or valuation, and a traversal is a sequence of
admissible coordinate updates. Voxel language fits the same role at nonary
scale: a local \(3^2\) or \(3^k\) block can be read as a small cell of
structured state, and overflow means that the current cell cannot preserve the
relevant invariant without adding a refinement coordinate.

Early implementations tested traversal continuity using sampled video/state
traces. A video trace is a sampled path through state space: continuity testing
asks whether adjacent frames preserve the expected local carrier relation,
whether an apparent jump is an admissible carry/refinement, and whether the
projection defect increases at the boundary. Such tests motivated carrier
choices and exposed discontinuities. The experiments are heuristic unless
bound to a reproducible metric and verification surface.

The sibling inventories record this boundary. Early implementations explored
ultrametric refinement, traversal continuity, diffusion rollout, branch-density
localization, and adaptive refinement. One tree-diffusion diagnostic records
\texttt{tree\_rollout\_mse = 4.246448141002339e-06} versus
\texttt{rbf\_rollout\_mse = 2.155684319973986e-05}. These experiments motivated
the refinement and transport structures later formalized in the carrier; they
are not theorem-level evidence.

The selected figures are copied into
\texttt{Docs/Images/paper1-origin-figures/} and recorded in
\texttt{Docs/papers/drafts/PaperDraftWorkingFolder/FigureCandidateManifest.md}. They
illustrate the refinement and traversal geometries that motivated the carrier
formalism.

\begin{center}
\centering
\begin{minipage}{0.48\linewidth}
\includegraphics[width=\linewidth]{Docs/papers/drafts/PaperDraftWorkingFolder/figures/paper1-origin-figures/dashifine_ultrametric_tree_frame0.png}
\end{minipage}\hfill
\begin{minipage}{0.48\linewidth}
\includegraphics[width=\linewidth]{Docs/papers/drafts/PaperDraftWorkingFolder/figures/paper1-origin-figures/dashifine_ultrametric_grid_frame0.png}
\end{minipage}
\par\smallskip\small\textbf{Figure A.1. Ultrametric refinement geometry.}
States sharing longer refinement
prefixes cluster naturally under the induced ultrametric; divergence occurs at
the first incompatible address or valuation depth.
\end{center}

\begin{center}
\centering
\begin{minipage}{0.48\linewidth}
\includegraphics[width=\linewidth]{Docs/papers/drafts/PaperDraftWorkingFolder/figures/paper1-origin-figures/dashitest_tree_diffusion_rollout_mse.png}
\end{minipage}\hfill
\begin{minipage}{0.48\linewidth}
\includegraphics[width=\linewidth]{Docs/papers/drafts/PaperDraftWorkingFolder/figures/paper1-origin-figures/dashitest_tree_diffusion_tree_band_quotient.png}
\end{minipage}
\par\smallskip\small\textbf{Figure A.2. Traversal continuity and tree-diffusion rollout diagnostics.}
The
sibling run gives a concrete visual and metric surface for tree-structured
transport; it does not promote a video-continuity or codec-correctness theorem.
\end{center}

\begin{center}
\centering
\begin{minipage}{0.48\linewidth}
\includegraphics[width=\linewidth]{Docs/papers/drafts/PaperDraftWorkingFolder/figures/paper1-origin-figures/fracdash_rank4_basin_branch_density.png}
\end{minipage}\hfill
\begin{minipage}{0.48\linewidth}
\includegraphics[width=\linewidth]{Docs/papers/drafts/PaperDraftWorkingFolder/figures/paper1-origin-figures/dashirtx_refinement_frame_100.png}
\end{minipage}
\par\smallskip\small\textbf{Figure A.3. Branch density and adaptive refinement.}
The diagnostics show
refinement concentrating around structured local regions rather than spreading
uniformly across the carrier.
\end{center}

UFT/trie refinement gives this traversal a compression geometry. A trie node is
a coarse address; descending the trie adds precision. Nearby branches share
prefix structure, and separation occurs at the first incompatible depth. This
matches the ternary ultrametric reading already used in the carrier spine:
refinement depth is not just a storage index but a metric and semantic
relation. A motif is then a reusable local shape in the trie, not an equality
claim.

\texttt{FactorVec} is the corresponding prime-coordinate refinement of the
same idea. Instead of only trit positions, a state carries valuation data
across tracked prime lanes. Multi-lane transport, chamber legality, signed
scans, and defect summaries turn traversal into a controlled carrier geometry:
a move may be legal or illegal, stable or repatterning, contractive or
expansive, and these labels are computed at the representation layer rather
than inferred from prose.

The transition from trits to carrier geometry is therefore historical and
mathematical. Ternary traversal made unresolvedness explicit; \texttt{Base369},
cube/voxel traversal, video-like continuity traces, UFT tries, and
\texttt{FactorVec} lanes added progressively sharper ways to represent
transition, address, refinement, and defect. These origins do not prove full
physics unification, automatic convergence, Monster representation semantics,
empirical adequacy, or domain recovery. They record the path by which stronger
claims would have to be typed, projected, tested, and closed.

\section{Conclusions}

DASHI is a constructive physics-unification program centered on a coherent
carrier geometry. Prime-coordinate structure, ultrametric refinement,
projection-defect decomposition, filtered operator surfaces, tracked lane
actions, local root-geometry boundaries, semantic compression, and bounded
empirical receipts are treated as parts of one mathematical construction.

The current frontier is substantial but not final. G2 records a
direction-indexed carrier schema, G3 records selected finite-support
subtraction support, G6 records tracked cross-lane commutation, and E8/LILA
records local root-boundary structure. UFT and motif structures keep refinement,
semantic distance, and residual structure in the same typed vocabulary; the
Drell-Yan receipt records bounded empirical contact and a stricter residual
shape law that still fails.

The residual obligations remain explicit throughout the construction, so
stronger physical claims require the corresponding adapters before promotion.
Paper 1's contribution is a constructive ultrametric carrier geometry whose
refinement, transport, projection, compression, and quotient structures support
staged physics-facing derivation surfaces with a visible certification
frontier.

\nocite{cms-smp-20-003,hepdata-ins2079374,collins-soper-sterman-1985,ebert-michel-stewart-tackmann-2021,drell-yan-1970,agda-stdlib,wellen-cartan-modal-hott-2018}
\bibliographystyle{plain}
\bibliography{Docs/papers/drafts/PaperDraftWorkingFolder/Paper1_References}

\end{document}
